%======================================================================
\section{A polynomial model}\label{sec:toy}\label{sec:model}
%======================================================================

In this section and the next three the Bessel kernel of the true family is replaced by $e^{-y}$. The resulting
\emph{model} is the instance $\varphi=e^{-y}$ of Definition~\ref{def:nested}(v): a Hermite interpolation problem for
polynomials, with the same nodes $y_j=2\pi n_j$. At leading order the true family \emph{is} such a problem
(Proposition~\ref{prop:dictionary}), and the kernel-free calculus of Sections~\ref{sec:nested} and~\ref{sec:chain} applies to
both. In the model, polynomial structure gives proofs that are not available for the true kernel: the front is proved for
every $J\le 80$ (Theorem~\ref{thm:PP80}, computer-assisted), and its arithmetic input is visible exactly, through prefix
sums of prime powers (Lemma~\ref{lem:trace}). Section~\ref{sec:model-tail} proves far-node asymptotics and an obstruction
to proving the tail by sign regularity alone, and Section~\ref{sec:W} shows that a natural sign property of the derivative
weights fails at the critical node density.

Nothing in these sections implies a statement about the true family; the relation between the two is the subject of
Section~\ref{ssec:dictionary} and Numerical observation~\ref{obs:toy-true}, and the quantitative agreement reported there is
numerical.

%----------------------------------------------------------------------
\subsection{Definition, umbral form and a radial Fourier interpretation}\label{ssec:toy-def}
%----------------------------------------------------------------------

Let $T_m(x)=\sum_{i=0}^m S(m,i)x^i$ be the Touchard polynomials \cite{Tou39,Com74}, where $S(m,i)$ are the Stirling
numbers of the second kind. Since $(x\frac{d}{dx})^m e^{x}=T_m(x)e^{x}$ and the Euler operator $\Eul=y\frac{d}{dy}$ is
invariant under $y\mapsto -y$, we have $\Eul^{m}e^{-y}=T_m(-y)e^{-y}$, and so
\begin{equation}\label{eq:tau}
  (-\Eul^2)^k e^{-y}=\tau_k(y)\,e^{-y},\qquad \tau_k(y):=(-1)^kT_{2k}(-y)=(-1)^k\Bigl[y^{2k}-\tbinom{2k}{2}y^{2k-1}+\cdots\Bigr],
\end{equation}
because $S(m,m-1)=\binom m2$. Thus $\tau_k$ has degree $2k$, leading coefficient $(-1)^k$ and $\tau_k(0)=0$ for
$k\ge1$; for instance $\tau_1=y-y^2$ and $\tau_2=y^4-6y^3+7y^2-y$.

\begin{definition}[The model]\label{def:toy}
\textup{(a)} For a polynomial $P(u)=\sum_{k=0}^d p_ku^k$ put
$Q_P(y):=e^{y}P(-\Eul^2)e^{-y}=\sum_k p_k\tau_k(y)$, a polynomial of degree $2\deg P$ with $Q_P(0)=P(0)$. In the
notation of Definition~\ref{def:nested} this is the family with base kernel $\varphi=e^{-y}$ on $I=(0,\infty)$, and
$\Fc_P=Q_Pe^{-y}$.

\textup{(b)} Let $\yv=(y_1<\dots<y_J)\subset(0,\infty)$ satisfy \hyp{N}$_{\yv}$, i.e.\ the matrix
$[\tau_k(y_j);\tau_k'(y_j)]_{j\le J,\,1\le k\le 2J}$ is non-singular. The \emph{model solution} is the unique
$P_{\yv}=1+\sum_{k=1}^{2J}p_k(\yv)u^k$ for which $Q_{\yv}:=Q_{P_{\yv}}$ has a double zero at every $y_j$. Put
$A_{\yv}(y):=\prod_j(y-y_j)^2$. If $p_{2J}(\yv)\neq0$, then $\deg Q_{\yv}=4J$, and the \emph{cofactor}
$R_{\yv}:=Q_{\yv}/A_{\yv}$ is a polynomial of degree $2J$ with leading coefficient $p_{2J}$ and
$R_{\yv}(0)=\prod_jy_j^{-2}$; we write $\hat R_{\yv}:=R_{\yv}/R_{\yv}(0)$ and $\tilde R_{\yv}:=R_{\yv}/p_{2J}$.

\textup{(c)} \Rp$_{\yv}$ means $R_{\yv}\in\RR_{>0}[y]$ (all $2J+1$ coefficients positive, in particular
$p_{2J}>0$); the model front \FRt$_{\yv}$ means $Q_{\yv}>0$ on $(0,\infty)\setminus\{y_1,\dots,y_J\}$. Clearly
\Rp$_{\yv}$ implies \FRt$_{\yv}$.

\textup{(d)} For an increasing sequence $y_1<y_2<\cdots$ write $\yle K:=(y_1,\dots,y_K)$ and $P_K,Q_K,R_K,p^{(K)}_k$
for its objects. \emph{$\TEL_J$} holds if $\hat R_J-\hat R_{J-1}$ has non-negative coefficients in degrees
$1,\dots,2J-2$, and \emph{strict $\TEL_J$} if it has positive coefficients in degrees $1,\dots,2J$.

\textup{(e)} \emph{$\PT_K$} is the condition $\sum_{j\le K}y_j>K(4K-1)$, and \emph{$\PT_\delta$ at $K$} is
$\sum_{j\le K}y_j\ge(1+\delta)K(4K-1)$. The progression $y^*_K:=\frac{21}{20}(8K-5)$ has
$\sum_{j\le K}y^*_j=1.05\,K(4K-1)$ for every $K$.

\textup{(f)} On the prime powers, $y_j=2\pi n_j$, $\yPP_J=(2\pi n_1,\dots,2\pi n_J)$, and $P_J,Q_J,R_J$ denote the
objects of $\yPP_J$.
\end{definition}

The degree, leading coefficient and constant term in (b) are read off from \eqref{eq:tau}: only $p_{2J}\tau_{2J}$
contributes to degree $4J$, and $Q_{\yv}(0)=P_{\yv}(0)=1$.

\begin{proposition}[Poisson, Mellin and umbral forms]\label{prop:umbral}
Let $P$ be a polynomial and $F_P:=Q_Pe^{-y}$.
\begin{enumerate}[label=\textup{(\alph*)}]
\item $F_P(y)=\sum_{j\ge0}P(-j^2)\,(-y)^j/j!$.
\item For $\Re s>0$, $\Gamma(s)^{-1}\int_0^\infty y^{s-1}F_P(y)\dd y=P(-s^2)$. Conversely, a polynomial
$Q=\sum_mq_my^m$ of degree $\le 4J$ equals $Q_P$ for some $P$ of degree $\le 2J$ if and only if
$M_Q(s):=\sum_mq_m\,s(s+1)\cdots(s+m-1)$ is an even polynomial.
\item Equivalently, $Q$ is of the form $Q_P$ if and only if, for $l=1,\dots,2J$,
$q_{2l-1}=-\sum_{k=2l}^{4J}|s(k,2l-1)|\,q_k$, where $s(k,i)$ are the Stirling numbers of the first kind. For
$l=2J$ this reads $q_{4J-1}=-\binom{4J}{2}q_{4J}$.
\item If $P(-j^2)=\sum_mb_m(j)_m$ with falling factorials $(j)_m$, then $Q_P(y)=\sum_mb_m(-y)^m$.
\end{enumerate}
\end{proposition}

The proof is in Appendix~\ref{app:proofs-toy}.

Let $\Hank g(y):=\int_0^\infty J_0(2\sqrt{yw})\,g(w)\dd w$, where $J_0$ is the Bessel function of order~$0$, and
let $\Lag$ be the linear map $y^m\mapsto m!\,L_m(y)$ on polynomials ($L_m$ the Laguerre polynomials).

\begin{proposition}[A two-dimensional radial Fourier form]\label{prop:radial}
Let $Q$ be a polynomial with $Q(0)=1$ and $F:=Qe^{-y}$. The following are equivalent:
\textup{(i)} $Q=Q_P$ for some polynomial $P$; \textup{(ii)} $\Hank F=-F'$; \textup{(iii)} $\Lag Q=Q-Q'$.
For $v\in\RR^2$ put $f(v):=F(\pi|v|^2)$ and $\FT f(\eta):=\int_{\RR^2}f(v)e^{-2\pi i\langle v,\eta\rangle}\dd v$. Then
$\FT f(\eta)=(\Hank F)(\pi|\eta|^2)$, so \textup{(i)} says $\FT f(\eta)=-F'(\pi|\eta|^2)$, and a double zero of $Q$ at
$y=2\pi n$ is equivalent to the vanishing of both $f$ and $\FT f$ on the circle $|v|=\sqrt{2n}$.
\end{proposition}

The proof is in Appendix~\ref{app:proofs-toy}.

So the model at the prime-power nodes asks for a radial function $f$ on $\RR^2$ with $\FT f=-F'(\pi|\cdot|^2)$ such
that $f$ and $\FT f$ vanish on the circles $|v|=\sqrt{2n}$, $n\in\PP_J$. For \emph{all} radii $\sqrt n$, $n\ge1$, Stoller
\cite{Sto21} proved that every Schwartz function on $\RR^d$, $d\ge2$, is determined by its restrictions, and those of its
Fourier transform, to the spheres of these radii, extending the one-dimensional interpolation theorem of Radchenko and
Viazovska \cite[Theorems~1 and~2]{RV19}; with derivatives at the radii $\sqrt{2n}$ the analogous statement is proved in
dimensions $8$ and $24$ \cite[Theorem~1.7]{CKMRV22}, and the general case of derivatives of order $<k$ at the radii $\sqrt{kn}$
is \cite[Open Problem~7.1]{CKMRV22}. Here the radii are $\sqrt{2n}$ for the prime powers $n$ only. Heuristically, all
integers $n$ form a node set of critical density for such problems, while the prime powers are sparse; the data
of Section~\ref{sec:W} are consistent with this picture, but we use it only as motivation. Expansions in Laguerre
polynomials of this kind, with prescribed double roots, are also the standard ansatz for sign-uncertainty problems
\cite{CG19}.

\begin{lemma}[The $u$-coefficient as an integral]\label{lem:p1-integral}
Let $P$ be a polynomial with $P(0)=1$ and $F=Q_Pe^{-y}$. Then $\int_0^\infty(F(y)-e^{-y})\frac{\dd y}{y}=0$ and
$p_1=-\int_0^\infty(F(y)-e^{-y})\log y\,\frac{\dd y}{y}$.
\end{lemma}

\begin{proof}
By Proposition~\ref{prop:umbral}(b), $\Gamma(s)(P(-s^2)-1)=\int_0^\infty(F(y)-e^{-y})y^{s-1}\dd y=:I(s)$ for
$\Re s>0$; since $F(y)-e^{-y}=(Q_P(y)-1)e^{-y}=O(y)$ at $0$, $I$ is analytic for $\Re s>-1$. As
$P(-s^2)-1=-p_1s^2+O(s^4)$ and $\Gamma(s)=s^{-1}+O(1)$, $I(s)=-p_1s+O(s^2)$; so $I(0)=0$ and $I'(0)=-p_1$.
\end{proof}

Thus $p_1$, the scalar that controls the tail in Section~\ref{ssec:toy-tail}, is a logarithmic moment of
$F-e^{-y}$.

%----------------------------------------------------------------------
\subsection{Dictionary to the true kernel}\label{ssec:dictionary}
%----------------------------------------------------------------------

Recall $\kt(z)=\sqrt{2z/\pi}\,e^{z}K_0(z)$ ($z>0$) from Section~\ref{ssec:farfield}, and from Fact~\ref{thm:bessel} that
$\Psi(\xi)=2\pi\sqrt x\sum_{N\ge1}d(N)\Wnull(2\pi Nx)$ with $\Wnull=\Eul^2(\Eul+1)^2K_0$, $x=e^{2\pi\xi}$.

\begin{proposition}[Toy--true dictionary]\label{prop:dictionary}
\textup{(a)} $\kt(z)=\pi^{-1/2}\int_0^\infty e^{-v}v^{-1/2}\bigl(1+\tfrac{v}{2z}\bigr)^{-1/2}\dd v$; hence $\kt$ is
increasing, $0<\kt<1$, and $\kt(z)=1-\frac1{8z}+\frac9{128z^2}+O(z^{-3})$ as $z\to\infty$.

\textup{(b)} For a polynomial $P$, with $P^{\Xi}(u)=P(u)(u+\frac14)^2$ as in Fact~\ref{lem:abel}, put
\[
  \mathfrak F_P(z):=P^{\Xi}(-\Eul^2)\bigl[e^{-z}\kt(z)\bigr].
\]
Then, with $y=2\pi x$,
\begin{equation}\label{eq:dictionary}
  \FT{\Xi^2P(t^2)}(\xi)=\pi\sum_{N\ge1}d(N)N^{-1/2}\,\mathfrak F_P(Ny).
\end{equation}
\textup{(c)} $\mathfrak F_P=\mathfrak m\bigl[Q_{P^{\Xi}}(z)e^{-z}\bigr]$, where $\mathfrak m$ is the Mellin multiplier
with symbol
\[
  m(s)=\frac{\Gamma(\frac{s+1/2}{2})^2}{\Gamma(\frac s2)\Gamma(\frac{s+1}2)};
\]
for $s>0$, $0<m(s)\le1$ and $m(s)\to1$ as $s\to\infty$.

In particular, if $\kt$ is replaced by $1$ and only the term $N=1$ is kept, the true transform becomes
$\pi\,Q_{P^{\Xi}}(y)e^{-y}$: the true family at leading order is the model for the polynomial $P^{\Xi}$, which has
a frozen double root at $u=-\frac14$.
\end{proposition}

The proof is in Appendix~\ref{app:proofs-toy}.

The dictionary does not make the true family a perturbation of the model in any useful norm: the two are different
interpolation problems, and a positive smoothing cannot transport positivity between them.

\begin{proposition}[No transfer by positive smoothing]\label{prop:smoothing-obstruction}
\textup{(a)} Let $k>0$ be measurable on $(1,\infty)$ and $T_kf(z):=\int_1^\infty f(zt)k(t)\frac{\dd t}{t}$ whenever
the integral converges absolutely. If $f$ is continuous, $f\ge0$ on $[z_0,\infty)$, $f\not\equiv0$ there, and
$T_kf(z_0)$ converges absolutely, then $T_kf(z_0)>0$.

\textup{(b)} $e^{-z}\kt(z)=T_kg(z)$ with $g(w):=w^{1/2}e^{-w}$ and
$k(t):=\pi^{-1/2}t^{1/2}(t-1)^{-1/2}\bigl(\frac{t+1}2\bigr)^{-1/2}>0$. Consequently, for every polynomial $P$,
\[
  \mathfrak F_P=T_k\bigl[w^{1/2}Q^{\mathrm{Ab}}_P(w)e^{-w}\bigr],
\]
where $Q^{\mathrm{Ab}}_P$ is the Abel preimage of Fact~\ref{lem:abel}.

\textup{(c)} Hence, if $P\neq0$ and $\mathfrak F_P(z_0)=0$ for some $z_0>0$, then $Q^{\mathrm{Ab}}_P$ is negative
somewhere on $(z_0,\infty)$.
\end{proposition}

The proof is in Appendix~\ref{app:proofs-toy}.

By \eqref{eq:dictionary} with $N=1$, $\FT{\Xi_\infty^2P(t^2)}(\xi)=\pi\mathfrak F_P(2\pi x)$. So for the Gamma-only family,
whose transform has double zeros at the nodes, the Abel preimage must take negative values beyond every node: the
``dips'' of the Abel preimage are forced. For the true kernel the transform at a node $y_n=2\pi n$ equals
$-\pi\sum_{N\ge2}d(N)N^{-1/2}\mathfrak F_P(Ny_n)$ rather than $0$, and these terms are of size $e^{-2y_n}$ up to
polynomial factors; we do not formalise the corresponding perturbation statement. In particular, no statement of
the form ``positivity in the model implies positivity for the true kernel'' can follow from positivity of the
smoothing alone, and the families are different solutions of different problems (in general the true $P_J$ is not
a model solution).

%----------------------------------------------------------------------
\subsection{The trace identity and small cases}\label{ssec:trace}
%----------------------------------------------------------------------

The node set enters the top of the cofactor only through the sum of the nodes.

\begin{lemma}[Trace identities]\label{lem:trace}
\textup{(a)} Let $\yv=(y_1,\dots,y_J)$ satisfy \hyp{N}$_{\yv}$ and $p_{2J}(\yv)\neq0$, and write
$R_{\yv}=p_{2J}(y^{2J}+\rho_1y^{2J-1}+\cdots)$. Then $\rho_1=2\sum_{j\le J}y_j-2J(4J-1)$; equivalently, the $2J$ roots
of $R_{\yv}$ sum to $2J(4J-1)-2\sum_jy_j$.

\textup{(b)} Consequently \Rp$_{\yv}$ implies $\PT_J$, and if \Rp\ holds for every prefix $\yle K$, $K\le J$, then
$\PT_K$ holds for every $K\le J$.

\textup{(c)} \textup{(chains)} For a regular chain $Z$ of level $K$ with monic cofactor $\pf_Z$ and $0$-augmented
cofactor $\pg_Z$ \textup{(}Definition~\ref{def:chain-laws}\textup{)}, $[y^{K-1}]\pf_Z=\Sigma Z-K(2K-1)$ and
$[y^K]\pg_Z=\Sigma Z-(K+1)(2K+1)$, where $\Sigma Z$ is the sum of the points of $Z$ with multiplicity.

\textup{(d)} \textup{(integer nodes)} For $y_j=2\pi(j+1)$, $\PT_J$ holds if and only if $J\le 12$. Hence for these
nodes \Rp\ fails for every $J\ge13$ at which \hyp{N} holds.
\end{lemma}

\begin{proof}
(a) By \eqref{eq:tau}, only $p_{2J}\tau_{2J}$ contributes to the coefficients of $y^{4J}$ and $y^{4J-1}$ of $Q_{\yv}$, so
$Q_{\yv}=p_{2J}\bigl(y^{4J}-\binom{4J}2y^{4J-1}+\cdots\bigr)$, while $A_{\yv}=y^{2J}-2(\sum_jy_j)y^{2J-1}+\cdots$;
divide, and note $\binom{4J}{2}=2J(4J-1)$. (b) Under \Rp, $p_{2J}>0$ and $p_{2J}\rho_1>0$. (c) $\omega_Z\pf_Z$ is a monic
element of $\Asp^{\mathrm{pol}}_K=\operatorname{span}(\tau_0,\dots,\tau_K)$ of degree $2K$, hence equals
$(-1)^K\tau_K$ plus lower $\tau$'s, i.e.\ $y^{2K}-K(2K-1)y^{2K-1}+\cdots$, and $\omega_Z=y^K-(\Sigma Z)y^{K-1}+\cdots$. The
second identity is the first for the chain $Z+0$ of level $K+1$. (d) $\sum_{j\le J}2\pi(j+1)=\pi J(J+3)$, and
$\pi J(J+3)>J(4J-1)$ if and only if $J<(3\pi+1)/(4-\pi)=12.14\ldots$; then apply (b).
\end{proof}

For the prime powers, $\sum_{j\le K}n_j\sim\frac12K^2\log K$, so the trace condition holds with a margin that grows
like $K^2\log K$; this is the only place where the arithmetic of the nodes enters the top coefficient. At $J=1$
everything is explicit. We use the property \textup{(W)} of Definition~\ref{def:W} below.

\begin{proposition}[One node]\label{prop:J1}
Let $J=1$, $y_1=a>0$ and $D(a):=2a^3-9a^2+12a-6$, which has a single real root $a_D=2.67765\ldots$. Then
\hyp{N} holds if and only if $D(a)\neq0$, and in that case
\begin{gather*}
  p_1=\frac{4a^3-18a^2+14a-1}{a^2D(a)},\qquad p_2=\frac{2a-1}{a^2D(a)},\qquad
  R=\frac{(2a-1)y^2+2(2a-1)(a-3)y+D(a)}{a^2D(a)},\\
  \mu_1=\frac{a^3-6a^2+7a-1}{aD(a)},\qquad Q''(a)=\frac{4(4a^3-12a^2+9a-3)}{a^2D(a)},
\end{gather*}
where $\mu_1$ is the derivative weight of the $p_1$-rule (Lemma~\ref{lem:dualrule}). Hence:
\textup{(a)} \Rp\ holds if and only if $a>3$, i.e.\ if and only if $\PT_1$;
\textup{(b)} $Q''(a)>0$ for every $a>a_D$;
\textup{(c)} for $a>a_D$, \textup{(W)} holds if and only if $a>4.49086\ldots$, the largest root of
$a^3-6a^2+7a-1=\tau_2(a)/a$;
\textup{(d)} the screening factor is
$S_{(a)}=(2a-1)(2a^7-33a^6+208a^5-620a^4+908a^3-636a^2+168a-24)/(a^2D(a)^2)$; it equals $1-8/a+8/a^2+O(a^{-3})$,
and for $a>a_D$ it is positive if and only if $a>3.26513\ldots$.
\end{proposition}

The proof (Appendix~\ref{app:proofs-toy}) is an explicit computation with $\tau_1$, $\tau_2$, $\tau_3$ and Cramer's rule.

So in the model \textup{(W)} is not a consequence of positivity: for $3<a<4.49$ the front holds but $\mu_1<0$.
Zero slack in the trace conditions is not enough, either.

\begin{remark}[Zero-slack prefix traces are not sufficient]\label{rem:PT0-false}
At $J=2$ the node sets $(y_1,y_2)=(3.02,11.3)$, $(3.05,11.05)$ and $(3.02,13.0)$ satisfy $\PT_1$ and $\PT_2$, but
the coefficient $[y]R_{\yv}$ is negative there (approximately $-5.0\cdot10^{-4}$, $-3.6\cdot10^{-4}$ and
$-1.7\cdot10^{-5}$; exact rational computation, \nolinkurl{toy/pt/verify_pt0_false.py} with
\nolinkurl{toy/pt/params/pt0_false.json}, which also certifies \hyp{N} at the three points). So $\PT_K$ for all $K\le J$ does not imply \Rp, and a relative slack
is needed. We use $\delta=0.05$.
\end{remark}

The three propositions below are exact certificates. Each one is a finite computation in exact integer arithmetic
whose design is described in Appendix~\ref{app:proofs-PT}; the positivity step is Lemma~\ref{lem:posc}(b).

\begin{proposition}[Positive cofactors for two nodes]\label{thm:PT2}
Let $J=2$ and $0<y_1<y_2$ with $y_1\ge\frac{63}{20}$ and $y_1+y_2\ge\frac{147}{10}$ \textup{(}that is,
$\PT_{0.05}$ at $K=1,2$\textup{)}. Then \hyp{N}$_{\yv}$ holds and $R_{\yv}\in\RR_{>0}[y]$; in particular $p_4>0$.
\end{proposition}

\begin{proposition}[Positive cofactors for three nodes]\label{thm:PT3}
Let $c=\frac{21}{20}$ and $0<y_1<y_2<y_3$ with $y_1\ge3c$, $y_1+y_2\ge14c$ and $y_1+y_2+y_3\ge33c$. Then
\hyp{N}$_{\yv}$ holds, $p_6>0$ and $R_{\yv}\in\RR_{>0}[y]$.
\end{proposition}

\begin{proposition}[Positive cofactors for four nodes]\label{thm:PT4}
Let $c=\frac{21}{20}$ and $0<y_1<y_2<y_3<y_4$ with $\sum_{j\le K}y_j\ge c\,K(4K-1)$ for $K=1,2,3,4$ \textup{(}i.e.\
$y_1\ge3.15$, $y_1+y_2\ge14.7$, $y_1+y_2+y_3\ge34.65$, $y_1+\dots+y_4\ge63$\textup{)}. Then \hyp{N}$_{\yv}$ holds,
$p_8>0$ and $R_{\yv}\in\RR_{>0}[y]$.
\end{proposition}

\begin{proof}[Proofs of Propositions~\ref{thm:PT2}--\ref{thm:PT4} \textup{(computer-assisted)}]
Fix $J\in\{2,3,4\}$ and $V:=\prod_{i<j}(y_i-y_j)$. The verifier computes, in exact arithmetic in
$\ZZ[y_1,\dots,y_J]$, polynomials $G_0,\dots,G_{2J}$ (the Cramer numerators of $D\cdot R_{\yv}$, divided by $V^4$) and
checks two identities: (i) $A_{\yv}(y)\sum_iG_iy^i$ satisfies the $2J$ odd Stirling conditions of
Proposition~\ref{prop:umbral}(c) identically; (ii) the determinant $D$ of the system in \hyp{N}$_{\yv}$ equals
$\epsilon_J\,(y_1\cdots y_J)^2V^4G_0$ with a sign $\epsilon_J=\pm1$. It then covers the region by finitely many rational
parametrisations of $[0,\infty)^m$ (two pieces for $J=2$, four for $J=3$, eleven for $J=4$; Appendix~\ref{app:proofs-PT})
and checks that every substituted, homogenised $G_i$ has only positive coefficients and a positive constant term; by
Lemma~\ref{lem:posc}(b), every $G_i>0$ on the region. Then $D\neq0$, i.e.\ \hyp{N}. By (i), $A_{\yv}\sum_iG_iy^i$ is an
admissible polynomial of degree $\le4J$ with double zeros at the nodes, hence a multiple of $Q_{\yv}$ (by \hyp{N} the
space of such polynomials is one-dimensional); comparing values at $0$ gives
$R_{\yv}=\sum_iG_iy^i/\bigl((y_1\cdots y_J)^2G_0\bigr)\in\RR_{>0}[y]$, and $p_{2J}=G_{2J}/((y_1\cdots y_J)^2G_0)>0$.

\emph{Ancillary data.} Script \nolinkurl{toy/pt/verify_pt.py} with parameter files \nolinkurl{toy/pt/params/pt2.json},
\nolinkurl{pt3.json}, \nolinkurl{pt4.json} (SHA-256 in Appendix~\ref{app:anc}); exact arithmetic (FLINT multivariate
polynomials over $\ZZ$); the sizes of the polynomials are in Table~\ref{tab:app-toycert}. For $J=2$ a second, independent route (Cramer's rule with explicit
ratios $r_i/r_4$, on the two original pieces) is \nolinkurl{toy/pt/verify_pt2_cramer.py}; for $J=4$ the region cover is
checked by \nolinkurl{toy/pt/check_pt4_region.py}. Re-run commands are listed in Appendix~\ref{app:anc}; the $J=4$
certificate takes about $40$ CPU-minutes.
\end{proof}

\begin{corollary}[The front for $J\le4$]\label{cor:PT-small}
For every $J\le4$ and every increasing $\yv$ satisfying $\PT_{0.05}$ at $K=1,\dots,J$, \hyp{N}$_{\yv}$, $p_{2J}>0$ and
\Rp$_{\yv}$ hold, hence \FRt$_{\yv}$. This applies to $\yPP_J$, $J\le4$, and to every increasing $\yv$ with
$y_k\ge y^*_k$ for all $k$.
\end{corollary}

\begin{proof}
For $J=1$, $y_1\ge3.15>a_D$ gives \hyp{N} and \Rp\ by Proposition~\ref{prop:J1}(a); for $J=2,3,4$ use
Propositions~\ref{thm:PT2}--\ref{thm:PT4}. The prime powers satisfy the four inequalities directly (their prefix sums $4\pi$, $10\pi$, $18\pi$,
$28\pi$ exceed $3.15$, $14.7$, $34.65$, $63$), and $y\ge y^*$ coordinatewise gives prefix sums $\ge1.05K(4K-1)$.
\end{proof}

The prime powers lie above the progression $y^*$, for every $K$.

\begin{proposition}[Prime powers dominate the $\PT$-tight progression]\label{prop:PP-dominates}
\textup{(a)} For every $K\ge1$, $2\pi n_K>y^*_K=\frac{21}{20}(8K-5)$; the smallest difference,
$2\pi n_7-y^*_7=18\pi-53.55>2.998$, occurs at $K=7$.
\textup{(b)} For every $K\ge1$, $\sum_{j\le K}2\pi n_j\ge1.2366\,K(4K-1)$; the ratio of the two sides is smallest at
$K=9$, where it equals $124\pi/315=1.23669\ldots$.
Consequently $\yPP$ satisfies $\PT_{0.05}$ \textup{(}indeed $\PT_{0.2366}$\textup{)} at every $K$.
\end{proposition}

The proof is in Appendix~\ref{app:proofs-toy}.

Theorem~\ref{thm:PP80} below settles the front on the prime powers for $J\le80$ by a different route. For general
node sets above $y^*$ there is a monotone-path reduction, which we record because it is the only route we know that
covers all such node sets.

\begin{remark}[A monotone-path reduction]\label{rem:Mss}
Let $\mathcal U_K:=\{\yv\in\RR^K:y_1<\dots<y_K,\ y_k\ge y^*_k\ \forall k\}$, regard node sets as doubled chains
(Definition~\ref{def:chain-laws}), and let $\yv\setminus y_j$ be $\yv$ with $y_j$ removed. Suppose that for every
$K\le J$:
$(\mathrm{AP}^+)_K$: $y^*_{\le K}$ is regular and $\tilde R(y^*_{\le K})\in\RR_{>0}[y]$;
$(\mathrm B)_K$: $T_K(y^*_K)>0$, where $T_K$ is the monic cofactor of the half-step chain
$(y^*_1,y^*_1,\dots,y^*_{K-1},y^*_{K-1},y^*_K)$;
$(\mathrm{LOC})_K$: $\partial\tilde R/\partial y_j\in\RR_{\ge0}[y]$ at every regular $\yv\in\mathcal U_K$ with
$\tilde R(\yv)\in\RR_{>0}[y]$ for which $\yv\setminus y_j$ is regular with $\tilde R(\yv\setminus y_j)\in\RR_{>0}[y]$.
Then every $\yv\in\mathcal U_J$ is regular with $\tilde R(\yv)\in\RR_{>0}[y]$, so \hyp{N}, $p_{2J}>0$ and \FRt\ hold
on $\mathcal U_J$.

The proof of this reduction is in Appendix~\ref{app:proofs-toy}.

The base data $(\mathrm{AP}^+)_K$ and $(\mathrm B)_K$ are certified for $K\le90$ (Proposition~\ref{prop:AP-B-cert}).
The local law $(\mathrm{LOC})$ is open; computed evidence is in Numerical observation~\ref{obs:misc}(a). For $J\le4$ the conclusion holds unconditionally by Corollary~\ref{cor:PT-small}.
\end{remark}

\begin{proposition}[Base data on the progression]\label{prop:AP-B-cert}
For every $K\le90$: the node set $y^*_{\le K}$, $y^*_k=\frac{21}{20}(8k-5)$, satisfies \hyp{N}, $p_{2K}>0$ and
$R\in\RR_{>0}[y]$ \textup{(}hence $(\mathrm{AP}^+)_K$\textup{)}; and the half-step cofactor $T_K$ at $y^*$ lies in
$\RR_{>0}[y]$, in particular $T_K(y^*_K)>0$ \textup{(}$(\mathrm B)_K$\textup{)}.
\end{proposition}

\begin{proof}[Proof \textup{(computer-assisted)}]
All nodes are rational, so everything is exact. For each $K\le90$, \nolinkurl{toy/ap/verify_ap.py} solves the
$2K\times2K$ system over $\QQ$ (FLINT \texttt{fmpq}) and checks the residual exactly, divides $Q$ by $A$ with zero
remainder, and checks $p_{2K}>0$ and the positivity of all $2K+1$ coefficients of $R$; as consistency checks it
verifies $r_1/r_0=1-P(-1)+2\sum_k1/y^*_k$ and $\rho_1=2\delta K(4K-1)$ exactly. \hyp{N} together with $p_{2K}\neq0$ gives
regularity, and $\tilde R=R/p_{2K}$. The script \nolinkurl{toy/ap/verify_halfstep.py} computes $T_K$ exactly in the
$\tau$-basis and checks its coefficients and $T_K(y^*_K)$. Parameter file \nolinkurl{toy/ap/params/ap_delta005.json}; logs
and re-run commands in Appendix~\ref{app:anc}.
\end{proof}

